\section*{Capítulo \ref{ch:b1:s-block} --- O hidrogênio e o bloco s}

\begin{solution}{exo:b1:s-block:1}
\ce{KH} e \ce{CaH2}: salinos; \ce{SiH4} e \ce{H2S}: covalentes;
\ce{PdH_{0.6}}: metálico. \ce{CaH2 + 2H2O -> Ca(OH)2 + 2H2}.
\end{solution}

\begin{solution}{exo:b1:s-block:2}
\ce{2Li + 2H2O -> 2LiOH + H2}, e o mesmo com Na e K;
\ce{Mg + 2HCl -> MgCl2 + H2}.
\end{solution}

\begin{solution}{exo:b1:s-block:3}
Básicos: \ce{Na2O}, \ce{MgO} (\ce{MgO + 2HCl -> MgCl2 + H2O}). Ácidos:
\ce{SiO2}, \ce{SO3}, \ce{CO2} (\ce{SO3 + H2O -> H2SO4}). Anfóteros:
\ce{Al2O3}, \ce{ZnO} (\ce{ZnO + 2OH- + H2O -> Zn(OH)4^2-}).
\end{solution}

\begin{solution}{exo:b1:s-block:4}
A pH 14, o par da água é \ce{2H2O + 2e- -> H2 + 2OH-}, $E = -0.83$~V;
$\log K = 2(-0.83 + 2.71)/0.059 = 64$.
\end{solution}

\begin{solution}{exo:b1:s-block:5}
Na chama, as colisões levam um elétron a um nível excitado; ao voltar ao
\omterm{def:b1:quantum-numbers:energy-level}{estado fundamental}, o átomo emite um fóton de energia igual ao
intervalo, específica do elemento. $E = hc/\lambda = 6.626 \times 10^{-34} \times
3.00 \times 10^8/(589 \times 10^{-9}) = \qty{3.375e-19}{J} = \qty{2.11}{eV}$.
% ledger: const:h, const:c, const:eV
\end{solution}

\begin{solution}{exo:b1:s-block:6}
$10^6/106.0 = \qty{9.43e3}{mol}$ de barrilha: \qty{9.43e3}{mol} de
calcário, \qty{0.944}{t}, e $2 \times 9.43 \times 10^3/0.90 =
\qty{2.10e4}{mol}$ de sal, \qty{1.23}{t}.
\end{solution}

\begin{solution}{exo:b1:s-block:7}
$[\ce{Mg^2+}] = 1.3/24.3 = \qty{0.053}{mol/L}$; $\mathrm{pH} = 14.00 -
\frac12(11.25 + \log 0.053) = 9.0$. \qty{53.5}{mol} de magnésio em
\qty{1.0}{m^3} exigem \qty{53.5}{mol} de \ce{Ca(OH)2}: \qty{4.0}{kg}.
\end{solution}

\begin{solution}{exo:b1:s-block:8}
A constante de equilíbrio (termodinâmica) é enorme para ambos; a
velocidade (cinética) não é decidida por $E^\circ$. O lítio funde a uma
temperatura mais alta, permanece sólido e é recoberto por uma camada de
produto menos solúvel, e perde seu elétron menos facilmente que o
potássio (energia de ionização maior): reage de modo regular, o potássio
violentamente.
\end{solution}

\begin{solution}{exo:b1:s-block:9}
$10^6/100.1 = \qty{9.99e3}{mol}$: CaO $9.99 \times 10^3 \times 56.1 =
\qty{0.560}{t}$; \ce{CO2} \qty{9.99e3}{mol}, $V = nRT/p = 9.99 \times 10^3
\times 8.314 \times 298.15/10^5 = \qty{248}{m^3}$.
\end{solution}

\begin{solution}{exo:b1:s-block:10}
$2.9 \times 10^8/6.9 = \qty{4.2e7}{kmol}$ de Li (massa em kg), $\times 73.8/2 =
\qty{1.55e9}{kg}$: cerca de 1.6 milhão de toneladas de carbonato de lítio.
$2.9 \times 10^8/8 = \num{3.6e7}$ baterias.
\end{solution}

\begin{solution}{exo:b1:s-block:11}
1.31, 0.84, 0.71 por cento a 20, 80, \qty{100}{\celsius}: a \omterm{def:b1:precipitation:solubility}{solubilidade}
diminui. A quente, menos carbonato fica em solução: recupera-se mais.
\end{solution}

\begin{solution}{exo:b1:s-block:12}
\ce{NaHCO3} é o menos solúvel dos sais que os íons podem formar, e a
solução fica saturada nele primeiro. A amônia torna a solução básica o
bastante para transformar \ce{CO2} em \ce{HCO3-} e ao mesmo tempo é
recuperável (como \ce{NH4Cl}, regenerado pela cal); o hidróxido de sódio
seria consumido e custaria mais que o produto.
\end{solution}

\begin{solution}{pb:b1:s-block:1}
\textbf{1.} Li \qty{6.0}{g/L}, Mg \qty{2.4}{g/L}.
\textbf{2.} Quatro metros cúbicos de salmoura dão um: três metros cúbicos
de água evaporam.
\textbf{3.} O sol e o clima seco fornecem a energia de graça.
\textbf{4.} O cloreto de sódio, de longe o sal mais abundante, satura
primeiro.
\textbf{5.} Li $6000/6.9 = \qty{870}{mol}$; Mg $2400/24.3 = \qty{98.8}{mol}$.
\textbf{6.} O magnésio precipitaria com o carbonato e contaminaria o
produto.
\textbf{7.} \ce{Mg^2+ + Ca(OH)2 -> Mg(OH)2 + Ca^2+}.
\textbf{8.} $98.8 \times 74.1 = \qty{7.32}{kg}$.
\textbf{9.} Restam $[\ce{Mg^2+}] = \num{9.88e-5}$ mol/L: $\mathrm{pH} = 14.00 -
\frac12(11.25 - 4.01) = 10.38$.
\textbf{10.} O hidróxido de lítio é solúvel: nenhum sólido se forma a esse pH.
\textbf{11.} \ce{Ca^2+ + CO3^2- -> CaCO3}.
\textbf{12.} $98.8 \times 106.0 = \qty{10.5}{kg}$.
\textbf{13.} \ce{2Li+ + CO3^2- -> Li2CO3}.
\textbf{14.} $435 \times 106.0 = \qty{46.1}{kg}$.
\textbf{15.} $435 \times 73.8 = \qty{32.1}{kg}$.
\textbf{16.} O carbonato de lítio é menos solúvel a quente.
\textbf{17.} $0.0084 \times 1000 = \qty{8.4}{kg}$.
\textbf{18.} Com pequenos volumes de água quente, em que ele é menos
solúvel.
\textbf{19.} $32.1 - 8.4 = \qty{23.7}{kg}$.
\textbf{20.} $23.7/32.1 = \qty{74}{\percent}$.
\textbf{21.} Ele ainda contém lítio: é devolvido às lagoas.
\textbf{22.} $2.9 \times 10^8/6.0 = \num{4.8e7}$ metros cúbicos.
\textbf{23.} \textbf{\qty{24}{kg}} (23.7) de carbonato de lítio por metro
cúbico.
\end{solution}
